\documentclass[UTF8,11pt]{ctexart}
\usepackage[a4paper,margin=2.2cm]{geometry}
\usepackage{hyperref}
\usepackage{enumitem}
\hypersetup{colorlinks=true,linkcolor=black,urlcolor=blue}
\setlength{\parskip}{0.5em}
\setlength{\parindent}{2em}
\title{《FlashSyn：基于反例驱动近似的闪电贷攻击合成》\\中文译读版}
\author{原文作者：Zhiyang Chen, Sidi Mohamed Beillahi, Fan Long}
\date{原文来源：ICSE 2024\\DOI：\url{https://doi.org/10.1145/3597503.3639190}}
\begin{document}
\maketitle

\section*{说明}
本文档依据用户提供的 PDF《FlashSyn: Flash Loan Attack Synthesis via Counter Example Driven Approximation》整理为中文译读版，重点翻译和重构论文的研究问题、核心技术、实验结果与结论，便于项目阅读与引用。为确保当前环境下可以稳定生成中文 PDF，文档采用单栏中文排版，不保留原始论文的版面、表格与公式细节，但保留关键术语与主要实验数字。

\section*{摘要}
在去中心化金融中，闪电贷允许借款人在单笔区块链交易内部临时借入大额资产，并在交易结束前偿还本金和费用。由于不需要预先存入抵押物，攻击者也可以利用闪电贷迅速聚集大量流动性，进而攻击存在设计缺陷的 DeFi 协议。本文提出 FlashSyn，这是一个面向闪电贷攻击的自动化攻击合成框架。它不是简单地检测某个单独合约中的编码错误，而是给定一组 DeFi 合约和候选动作后，自动搜索一个由多步函数调用组成的攻击序列，并同时求解每一步调用参数，使该攻击获得最大利润。

为了绕开 DeFi 协议逻辑过于复杂、难以直接用精确符号约束建模的问题，论文提出“基于近似的合成”方法：首先在分叉链环境中执行候选动作，采集输入输出数据点；然后用数值方法近似协议行为，包括多项式近似与最近邻插值；接着依据近似函数构造全局优化问题，搜索最有利可图的动作顺序与参数；最后再把合成出的攻击回放到真实协议执行环境中验证。如果估计利润与真实利润差异过大，就把该候选向量当作反例，继续收集新数据点修正近似。作者将这一框架实现为工具 FlashSyn，并在 16 个真实闪电贷攻击协议和 2 个 Damn Vulnerable DeFi 基准上进行了评估。结果表明，FlashSyn 可以在 18 个基准中的 16 个上自动合成出盈利攻击。

\section{引言}
论文从 DeFi 安全的一个核心难点出发：许多闪电贷攻击并不是传统意义上的“合约实现 bug”，而是跨多个协议、依赖大量临时资金、利用经济或设计缺陷的组合型攻击。对这种攻击，仅仅分析单个合约是不够的。攻击过程往往从一个闪电贷动作开始，中间穿过多个 DEX、借贷池、预言机或资产管理协议，最后再归还贷款并留下利润。由于涉及的动作序列和参数空间都非常大，自动化发现这类攻击极具挑战。

作者指出，现有静态分析、符号执行和验证工具大多聚焦于实现错误，例如重入、整数溢出或访问控制，而很难处理闪电贷攻击这种“跨协议、参数驱动、以利润为目标”的设计缺陷。即便攻击动作序列已经知道，仅仅为了找到能盈利的一组参数，也可能需要面对非常复杂的数学函数和状态依赖关系，导致求解器超时或无法收敛。

因此，FlashSyn 的核心问题被定义为：在不知道准确攻击向量的前提下，如何从候选动作集合中自动合成一个动作序列，并求出每一步的最佳参数，使整个交易在现实执行中获得正利润。

\section{背景与动机示例}
论文首先回顾了区块链、智能合约与 DeFi 的基础概念，然后用 Harvest Finance 的 USDC 金库攻击作为动机示例。该攻击发生在 2020 年 10 月 26 日，攻击者利用 Uniswap 提供的闪电贷能力，借入巨量 USDT 和 USDC，再与 Curve Y Pool、Harvest 金库等协议交互，操纵资产估值，最终获利约 30.7 万美元。

作者借这个案例说明三类核心难点。第一，\textbf{交互复杂}：攻击不是调用一个函数就结束，而是由一串相互影响的动作组成，前面一步修改的状态会影响后一步的利润。第二，\textbf{闭源组件}：某些被调用的外部合约并没有完整公开源码，使精确建模更加困难。第三，\textbf{数学复杂度高}：DeFi 协议中经常包含复杂的数值计算，例如 Curve 稳定币池的不变量计算需要求解高次方程，直接用符号执行工具处理会失败。论文报告说，即使对简化后的数学函数应用智能合约符号执行工具，也仍然会触发与求解器或内存相关的失败。

这些观察共同支持了 FlashSyn 的技术路线：不要强求对协议行为做完全精确的符号建模，而是用足够准确、计算更轻的近似模型来驱动攻击合成。

\section{问题形式化}
论文使用带标签迁移系统来抽象智能合约行为，把每一个用户可调用的方法都看作一个操作标签。给定某个链上状态和一组候选动作，工具需要寻找一个“动作向量”，即一个按顺序排列的函数调用序列，并为其中的符号参数赋上具体整数值，使系统从初始状态转移到新状态后，攻击者地址持有的资产总价值高于初始状态。

在这个形式化下，攻击目标可以表示为一个优化问题：最大化攻击者在终态与初态之间的余额差值，同时满足所有必要约束，例如任意地址的 token 余额不能为负、某些调用前必须先持有特定 share token，否则交易会回滚。这个抽象把“寻找攻击”从纯粹的程序路径搜索，转换为“动作序列枚举 + 约束优化 + 真实执行校验”的组合问题。

\section{FlashSyn 方法}
\subsection{整体流程}
FlashSyn 的主流程可以概括为五步。第一步，对每个候选动作收集一批初始数据点。每个数据点由动作执行前的关键状态、输入参数、执行后的状态变化和输出构成。第二步，根据这些数据点近似每个动作的状态转移函数。第三步，枚举一定长度以内的动作序列，并使用启发式规则剪枝。第四步，把近似函数装配成优化问题，用全局优化器搜索参数。第五步，将合成出的攻击向量放到分叉链中真实执行，验证其估计利润是否接近真实利润；如果偏差过大，则生成反例并进入下一轮修正。

\subsection{基于近似的合成}
论文最重要的创新是 synthesis-via-approximation，也就是“先近似、再合成”。对每个动作，FlashSyn 从实际执行收集数据点，再用两种数值方法学习动作效果：一是基于线性回归的多项式近似，二是最近邻插值。近似对象包括函数参数与关键状态变量对输出 token 数量、余额变化、后续状态的影响。

这种做法的关键好处是：优化器处理的是近似后的显式数值表达式，而不是直接面对复杂 EVM 字节码或难以求解的符号表达式。论文显示，这种近似足以驱动搜索过程，而且在很多案例中能得到与真实攻击接近、甚至更高利润的向量。

\subsection{动作序列枚举与剪枝}
FlashSyn 不是盲目穷举所有动作排列，而是在枚举动作向量时应用了多条启发式规则。第一，禁止相邻的重复动作，因为很多情况下两次连续调用同一方法可以合并成一次更大的调用。第二，限制同一动作在向量中的最大使用次数，避免无意义地拉长搜索空间。第三，检查必要前置条件，例如某些提现动作要求用户先持有 share token，那么序列中就必须先出现铸造该 share token 的动作。论文报告说，这些启发式能显著压缩搜索空间。

\subsection{反例驱动近似修正}
仅凭初始数据点学到的近似模型并不总是足够准确。优化器有时会在训练样本之外的区域找到看起来利润很高、但真实执行后效果不佳的参数组合。为此，FlashSyn 引入反例驱动近似修正机制。若某个攻击向量在真实链上执行得到的利润与估计利润差异超过阈值，就把它视为 counterexample，并从出错的动作开始向前回溯，补充新的输入输出数据点，再重新拟合近似模型。这个过程会多轮迭代，直到收益不再改进或达到超时阈值。

作者把这种机制类比为“对近似模型做 CEGAR 风格修正”，但这里修正的不是抽象状态图，而是用于驱动优化搜索的数值近似函数。

\section{FlashFind：候选动作识别}
论文还实现了一个配套组件 FlashFind，用来帮助用户自动筛选可能参与攻击的候选动作。FlashFind 的目标不是直接给出完整攻击，而是缩小搜索空间。

其工作大致分三步。首先，读取目标合约的 ABI，过滤掉 \texttt{view}、\texttt{pure}、仅管理员可调用或只影响权限控制的函数。其次，从历史交易中提取函数级 trace，学习某些难以穷举的特殊参数，例如地址、枚举值、字节串等。最后，在本地分叉链环境执行这些动作，验证它们是否可执行，并基于存储读写关系分析动作之间的 Read-After-Write 依赖，从而识别哪些动作真正可能相互作用。输出结果是一组“只剩下整数参数待求解”的动作候选，为 FlashSyn 后续合成提供输入。

\section{实现}
FlashSyn 使用 Python 实现，整体由五个组件构成。\textbf{Approximator} 负责学习动作近似模型；\textbf{Optimizer} 把近似模型转化为全局优化问题，并调用 \texttt{scipy.optimize.shgo} 搜索参数；\textbf{Runner} 在分叉链上执行交易，既用于初始与反例数据采集，也用于验证优化结果；\textbf{Synthesizer} 负责动作序列枚举、剪枝、分数排序与迭代调度；可选组件 \textbf{FlashFind} 则用于候选动作识别。论文提到，Runner 基于 Foundry 实现，而近似方法则依赖 \texttt{sklearn} 与 \texttt{scipy}。

一个重要实现细节是优先级评分。FlashSyn 不会对所有枚举序列都投入同样多的优化预算，而是先用较弱的参数搜索快速筛掉明显无望的序列，再把更多计算资源分配给当前已经显示出正利润潜力的动作向量。这使得它在有限时间预算内能够覆盖更多可能的攻击结构。

\section{实验评估}
论文在 18 个基准上评估 FlashSyn，其中 16 个是真实发生过闪电贷攻击的协议案例，另有 2 个来自 Damn Vulnerable DeFi。历史攻击横跨 Ethereum、BSC 与 Fantom，涉及 Harvest、bZx、Eminence、ValueDeFi、Warp、Puppet 等著名案例。作者把这些攻击在本地分叉链环境中重放，作为 ground truth；同时实现了一个需要人工推导精确数学摘要的 baseline，用来与近似方法比较。

实验围绕四个研究问题展开。第一，FlashSyn 是否能有效合成闪电贷攻击。结果表明，采用多项式近似的 FlashSyn-poly 能在 18 个基准中解决 16 个，平均归一化利润达到 0.945；采用插值方法的 FlashSyn-inter 能解决 13 个，平均归一化利润为 0.641。第二，近似驱动方法与精确摘要基线相比表现如何。论文指出，很多案例要么包含闭源组件，要么逻辑过于复杂，精确基线无法建立；而在可以比较的场景中，FlashSyn 得到的最佳利润平均达到精确基线的 97\%，个别案例甚至超过基线。第三，反例驱动修正是否有效。结果显示，开启反例循环后，FlashSyn-poly 的平均归一化利润从 0.717 提升到 0.945，并且能解决更多基准。第四，FlashFind 与 FlashSyn 结合时是否仍然有效。结果显示，在 14 个可用历史交易的基准上，FlashFind 能为 11 个基准识别出合理数量的候选动作，而 FlashSyn 在这些候选集合上仍能合成出盈利攻击，某些情况下还找到了与原始攻击不同的新动作组合。

论文还提到，一个作者与 Quantstamp 合作将 FlashSyn 应用于真实审计工作，在 3 个月内发现了两个零日闪电贷漏洞。这一结果从工业角度强化了工具的实用价值。

\section{局限性与相关工作}
论文承认，FlashSyn 仍然面临典型的可扩展性挑战。随着候选动作数量增长、攻击向量长度变长，搜索空间会呈指数级扩张。作者建议在实际审计中以模块化方式聚焦于相互依赖的动作集合，从而把问题规模控制在可处理范围内。

在相关工作方面，作者把现有研究分为三类。第一类是人工参数优化方法，即人工提取函数数学关系，再交给优化器求值；这类方法依赖专家手工推导，扩展性差。第二类是静态分析与符号执行工具，例如 Slither、Manticore 等，它们更擅长发现单合约实现错误，不适合处理跨协议的闪电贷攻击。第三类是更一般的程序合成与攻击合成研究。FlashSyn 的创新点，在于首次把“数值近似 + 反例修正 + 真实链验证”组合成一个面向 DeFi 闪电贷攻击的端到端自动合成框架。

\section{结论}
这篇论文的核心贡献不是提出某个新的漏洞类型，而是提出一种新的分析路线：面对逻辑极其复杂、难以精确求解的 DeFi 协议，不必执着于完整符号化，而可以通过数值近似学习动作行为，用近似模型驱动攻击合成，再通过反例不断修正。实验结果说明，这条路线在真实世界的闪电贷攻击案例上具有很强效果。对于你们的课程项目而言，这篇论文最大的启发不是去直接复现攻击合成，而是理解 trace、动作级建模、候选动作筛选和跨协议交互分析在 DeFi 安全中的重要性。

\section*{引用信息}
原始论文：Zhiyang Chen, Sidi Mohamed Beillahi, and Fan Long. \textit{FlashSyn: Flash Loan Attack Synthesis via Counter Example Driven Approximation}. In 2024 IEEE/ACM 46th International Conference on Software Engineering (ICSE '24), 2024.

\end{document}
